\chapter{Curriculum aplicado de 12 semanas}

\section{Principio}


Este curriculum no existe para "aprender mas". Existe para mejorar empleabilidad. Cada semana tiene un skill, una practica y una evidencia. Si la evidencia no se produce, la semana no cuenta como completada.


\subsection{Semana 1: narrativa tecnica}


Objetivo: explicar tu perfil sin sonar disperso.


Tareas:


\begin{itemize}
\item Escribir pitch de 30, 60 y 120 segundos.
\item Traducir el pitch a headline, CV summary y portfolio hero.
\item Eliminar terminos que inflen seniority.
\item Practicar explicacion de TwinSight sin abrir la tesis.
\end{itemize}


Evidencia: pitch escrito, grabacion de prueba, headline definitivo.


\subsection{Semana 2: Unity WebGL constraints}


Objetivo: poder defender por que WebGL importa y que restricciones impone.


Tareas:


\begin{itemize}
\item Revisar build settings, compression, memory and browser constraints.
\item Documentar 5 decisiones tecnicas de TwinSight relacionadas con WebGL.
\item Tomar una captura de build/performance si existe.
\item Preparar respuesta: what makes WebGL harder than desktop Unity?
\end{itemize}


Evidencia: seccion en README o portfolio sobre WebGL constraints.


\subsection{Semana 3: CAD-to-realtime optimization}


Objetivo: convertir el pipeline en prueba de Technical Artist.


Tareas:


\begin{itemize}
\item Preparar before/after visual.
\item Describir cleanup, hierarchy, mesh/material decisions.
\item Explicar que se optimizo y que no.
\item Preparar respuesta: how would you productionize the pipeline?
\end{itemize}


Evidencia: breakdown visual + bullet tecnico en CV.


\subsection{Semana 4: interaction and inspection UX}


Objetivo: mostrar que no solo importaste un modelo 3D; construiste una experiencia.


Tareas:


\begin{itemize}
\item Grabar selection, exploded view, clipping and visual modes.
\item Explicar para que sirve cada feature.
\item Relacionar features con problema de assembly inspection.
\item Preparar respuesta: how did you choose these interactions?
\end{itemize}


Evidencia: 4 GIFs/videos cortos o clips en demo.


\subsection{Semana 5: UI Toolkit / interface clarity}


Objetivo: que el proyecto parezca herramienta tecnica, no escena de prueba.


Tareas:


\begin{itemize}
\item Pulir paneles visibles.
\item Reducir texto innecesario.
\item Asegurar estados claros: selected, hovered, active mode.
\item Preparar respuesta: how would you improve the UI for production?
\end{itemize}


Evidencia: screenshot antes/despues o nota de UI.


\subsection{Semana 6: README and technical documentation}


Objetivo: que GitHub venda sin entrevista.


Tareas:


\begin{itemize}
\item README con overview, demo, features, architecture, pipeline, limitations.
\item Agregar imagenes.
\item Agregar roadmap tecnico.
\item Revisar que no haya claims falsos.
\end{itemize}


Evidencia: README final v1.


\subsection{Semana 7: visual presentation / ArtStation}


Objetivo: parecer technical artist, no solo developer.


Tareas:


\begin{itemize}
\item Elegir 5-8 imagenes fuertes.
\item Ordenarlas como historia de pipeline.
\item Agregar Human CGCookie como soporte visual.
\item Publicar breakdown o dejarlo listo.
\end{itemize}


Evidencia: ArtStation post o PDF/imagenes preparadas.


\subsection{Semana 8: interview defense}


Objetivo: defender el proyecto bajo presion.


Tareas:


\begin{itemize}
\item Practicar 20 preguntas.
\item Grabar walkthrough de 5 minutos.
\item Preparar respuestas sobre AI, digital twin, seniority, salary, EU authorization.
\item Corregir respuestas largas.
\end{itemize}


Evidencia: answer bank personal.


\subsection{Semana 9: application writing}


Objetivo: adaptar sin perder tiempo.


Tareas:


\begin{itemize}
\item Crear 3 versiones de mensaje: recruiter, technical lead, founder.
\item Crear 3 bullets CV por tipo de rol.
\item Crear plantilla de follow-up.
\item Enviar al menos 5 aplicaciones.
\end{itemize}


Evidencia: aplicaciones reales registradas.


\subsection{Semana 10: test assignment readiness}


Objetivo: no entrar indefenso a pruebas tecnicas.


Tareas:


\begin{itemize}
\item Preparar mini escena o checklist de prueba Unity.
\item Practicar explicar tradeoffs.
\item Definir limites para pruebas no pagadas.
\item Preparar respuesta para pedir scope y criterios.
\end{itemize}


Evidencia: test readiness note.


\subsection{Semana 11: negotiation readiness}


Objetivo: separar deseo, piso y valor.


Tareas:


\begin{itemize}
\item Escribir piso: 3 millones COP si experiencia afin.
\item Escribir ideal: USD 2000.
\item Escribir ambicioso: USD 6000.
\item Preparar respuestas a low offer y contractor terms.
\end{itemize}


Evidencia: scorecard listo.


\subsection{Semana 12: pivot audit}


Objetivo: decidir con datos.


Tareas:


\begin{itemize}
\item Contar aplicaciones, respuestas, entrevistas, rechazos.
\item Separar por rol: Unity TA, WebGL, visualization, XR, Python.
\item Detectar que target responde mejor.
\item Elegir siguiente ciclo de 4 semanas.
\end{itemize}


Evidencia: decision escrita en tracker.
